\topic{1}{Konstrukcja i analiza portfela papierów wartościowych}{
Portfel opisują wagi $w_i$ poszczególnych aktywów,
zwykle spełniające $\sum_iw_i=1$.
Jego oczekiwana stopa zwrotu to
$\mu_p=\mathbf w^\top\boldsymbol\mu$,
a wariancja
$\sigma_p^2=\mathbf w^\top\Sigma\mathbf w$,
gdzie $\Sigma$ jest macierzą kowariancji
zwrotów. Korelacje decydują,
ile ryzyka można ograniczyć
przez dywersyfikację.

W ujęciu Markowitza wybiera się
portfele efektywne: dla danego
ryzyka mają najwyższy oczekiwany
zwrot albo dla danego zwrotu
najmniejsze ryzyko.
Analiza wymaga oszacowania
$\boldsymbol\mu$ i $\Sigma$
z danych, określenia horyzontu,
kosztów transakcyjnych
oraz ograniczeń wag.
Wynik może być niestabilny,
ponieważ błędy oszacowania
średnich zwrotów są duże;
warto sprawdzać odporność
rozwiązania na zmianę
okna danych.}
\topic{2}{Portfel najmniejszego ryzyka i portfel optymalny}{
Portfel minimum wariancji
rozwiązuje problem
$\min_{\mathbf w}\mathbf w^\top
\Sigma\mathbf w$ przy
$\mathbf 1^\top\mathbf w=1$.
Bez zakazu krótkiej sprzedaży
i przy odwracalnej $\Sigma$
rozwiązanie ma postać
$\mathbf w_{\rm GMV}
=\Sigma^{-1}\mathbf 1/
(\mathbf 1^\top\Sigma^{-1}\mathbf 1)$.

Portfel optymalny wymaga
dodatkowego kryterium
preferencji inwestora,
na przykład maksymalizacji
$\mathbf w^\top\boldsymbol\mu
-\lambda\mathbf w^\top\Sigma\mathbf w/2$.
Parametr $\lambda$ wyraża
awersję do ryzyka.
Portfel minimum wariancji
nie musi maksymalizować
tej funkcji. Dodatkowe
ograniczenia, koszty
i niedokładność estymacji
mogą znacznie zmienić
wagi, dlatego rozwiązanie
algebraiczne jest punktem
wyjścia do analizy,
a nie prognozą wyników.}
\newpage
\topic{3}{Korelacje, zwroty i taksonomia portfela}{
Z cen $P_i(t)$ oblicza się
zwroty proste lub logarytmiczne,
np. $r_i(t)=\ln P_i(t)-
\ln P_i(t-\Delta t)$.
Korelacja Pearsona
$\rho_{ij}$ mierzy liniową
współzmienność zwrotów,
nie samych poziomów cen.
Korelacje cen często bywają
pozorne z powodu trendów.

Odległość
$d_{ij}=\sqrt{2(1-\rho_{ij})}$
pozwala zbudować drzewo
minimalne lub hierarchię
podobieństw między aktywami.
Taka taksonomia może ujawnić
grupy sektorowe i wspólne
czynniki, ale zależy
od okna czasowego,
częstotliwości danych
i sposobu estymacji.
Korelacja nie mierzy
zależności przyczynowej
ani dobrze nie opisuje
wszystkich ekstremalnych
współruchów. Przed
porównaniem serii trzeba
zsynchronizować notowania
i obsłużyć brakujące dane.}
\topic{4}{Krachy jako analogia przejść fazowych}{
W ekonofizyce krach bywa
modelowany jako gwałtowna
zmiana stanu układu
złożonego z oddziałujących
uczestników. Analogią
przejścia fazowego są
narastające korelacje,
niestabilność i możliwe
zachowanie skalowe.
Modele logarytmiczno-
periodyczne próbują
opisać przyspieszanie
bańki przed punktem
krytycznym.

Nie ma jednak uniwersalnego
prawa fizycznego,
które wyznacza datę krachu.
Dopasowanie modelu
do historii może łatwo
przeuczyć dane,
a wybór początku bańki
zmienia wynik.
Rzetelne badanie wymaga
definicji krachu,
testów poza próbą,
przedziałów niepewności
i porównania z prostszymi
modelami. Wskaźnik
ostrzegawczy jest
hipotezą statystyczną,
nie pewną zapowiedzią
przyszłego zdarzenia.}
\newpage
\topic{5}{Statystyczne własności rozkładów zwrotów}{
Zwroty finansowe mają
często rozkład o grubych
ogonach względem rozkładu
normalnego. Duże odchylenia
występują częściej niż
przewiduje model Gaussa.
Znaki krótkookresowych
zwrotów mogą mieć słabą
autokorelację, podczas
gdy wartości bezwzględne
lub kwadraty zwrotów
wykazują klastrowanie
zmienności.

Analiza obejmuje średnią,
wariancję, skośność,
kurtozę i funkcję
autokorelacji, ale
estymatory zależą
od okresu i rozdzielczości
danych. Rozkłady zwrotów
dodatkowo zmieniają się
w czasie, więc jedna
funkcja gęstości nie musi
opisywać całego rynku.
Trzeba sprawdzić wpływ
ekstremów, świąt,
braku płynności
i zmian metodologii.
Ciężki ogon sam w sobie
nie dowodzi konkretnego
mechanizmu generowania
cen.}
\topic{6}{Modele AR, MA i ARMA a rozkład Gaussa}{
Model AR($p$) opisuje
bieżący zwrot przez
poprzednie zwroty
i innowację:
$r_t=c+\sum_{i=1}^p
\phi_i r_{t-i}+\varepsilon_t$.
Model MA($q$) używa
bieżącej i poprzednich
innowacji, a ARMA($p,q$)
łączy oba składniki.
Parametry opisują
strukturę warunkowej
średniej szeregu.

Jeśli innowacje są
niezależne i normalne,
stacjonarny liniowy
model ARMA ma
gaussowskie rozkłady
skończenie wymiarowe.
Sam zapis ARMA
nie wymusza jednak
normalności; można
użyć innowacji
o grubych ogonach.
Model dla średniej
nie wystarcza
do opisu zmiennej
w czasie wariancji.
Dobór rzędów
opiera się na
diagnostyce reszt,
informacji kryterialnej
i sprawdzeniu
stacjonarności.}
\newpage
\topic{7}{Modele ARCH i GARCH}{
ARCH zakłada, że
warunkowa wariancja
zależy od przeszłych
kwadratów innowacji,
np. $h_t=\omega+
\alpha\varepsilon_{t-1}^2$.
GARCH dodaje poprzednią
wariancję:
$h_t=\omega+
\alpha\varepsilon_{t-1}^2+
\beta h_{t-1}$.
Zapis
$\varepsilon_t=\sqrt{h_t}z_t$
oddziela zmienną
skalę od standaryzowanej
innowacji $z_t$.

Dla prostego GARCH(1,1)
przy $\omega>0$,
$\alpha,\beta\geq0$
i $\alpha+\beta<1$
istnieje skończona
wariancja bezwarunkowa
$\omega/(1-\alpha-\beta)$.
Model uchwytuje
klastrowanie zmienności,
ale klasyczna wersja
symetrycznie reaguje
na dodatnie i ujemne
szoki tej samej wielkości.
W praktyce sprawdza się
reszty standaryzowane,
stabilność parametrów
i jakość prognoz
poza próbą.}
\topic{8}{Skalowanie, fraktalność i multifraktalność}{
Skalowanie oznacza
systematyczną zależność
statystyk przy zmianie
skali czasu. Dla
prostego procesu
samopodobnego moment
rzędu $q$ może spełniać
$E|r_{\Delta t}|^q
\propto(\Delta t)^{\zeta(q)}$.
W monofraktalnym
przybliżeniu $\zeta(q)$
jest liniowe w $q$;
multifraktalność
wiąże się z
nieliniowością tej
funkcji.

Szeregi finansowe
mogą pokazywać
zależność zmienności
na wielu skalach.
Modele multifraktalne
próbują połączyć
grube ogony,
klastrowanie
i skalowanie.
Estymacja jest jednak
wrażliwa na długość
próby, zakres skal
i skrajne obserwacje.
Pozorna multifraktalność
może też wynikać
ze zmian reżimu.
Wnioski wymagają
porównania z
danymi losowo
przemieszanymi
i modelami
odniesienia.}
\newpage
\section*{Podstawowe źródła do części III}
\begin{itemize}
\item H. Markowitz,
\href{https://onlinelibrary.wiley.com/doi/10.1111/j.1540-6261.1952.tb01525.x}
{Portfolio Selection}, 1952.
\item R. N. Mantegna,
\href{https://arxiv.org/abs/cond-mat/9802256}
{Hierarchical Structure in Financial Markets}, 1998.
\item R. F. Engle,
\href{https://web.pdx.edu/~crkl/ec572/readings/Engle82.pdf}
{Autoregressive Conditional Heteroscedasticity}, 1982.
\item T. Bollerslev,
\href{https://www.sciencedirect.com/science/article/pii/0304407686900631}
{Generalized Autoregressive Conditional Heteroskedasticity}, 1986.
\item E. Bacry, J. Delour, J.-F. Muzy,
\href{https://arxiv.org/abs/cond-mat/0005405}
{A multifractal random walk}, 2000.
\end{itemize}
